% TikZ 图：碰撞模型
% Day 10：冲量、动量与动量守恒
\begin{tikzpicture}[scale=0.9,>=Latex]
  % === 碰撞前 ===
  % m1
  \draw[thick] (-3,0.5) rectangle (-1.5,1.7);
  \node at (-2.25,1.1) {$m_1$};
  \draw[->,thick,red] (-1.5,1.1) -- (-0.3,1.1) node[above] {$v_1$};
  % m2
  \draw[thick] (0.5,0.5) rectangle (2.0,1.7);
  \node at (1.25,1.1) {$m_2$};
  \draw[->,thick,red] (2.5,1.1) -- (3.5,1.1) node[above] {$v_2$};
  \node[below] at (-0.75,0.3) {碰撞前};
  
  % === 碰撞后 ===
  % m1'
  \draw[thick] (-2,-2.0) rectangle (-0.5,-0.8);
  \node at (-1.25,-1.4) {$m_1$};
  \draw[->,thick,blue] (-0.5,-1.4) -- (0.8,-1.4) node[above] {$v_1'$};
  % m2'
  \draw[thick] (1.5,-2.0) rectangle (3.0,-0.8);
  \node at (2.25,-1.4) {$m_2$};
  \draw[->,thick,blue] (3.0,-1.4) -- (4.2,-1.4) node[above] {$v_2'$};
  \node[below] at (0.75,-2.2) {碰撞后};
  
  % 动量守恒公式放在右侧空白
  \node[red] at (3.8,-2.8) {$m_1 v_1 + m_2 v_2 = m_1 v_1' + m_2 v_2'$};
\end{tikzpicture}
